\backmatter
\chapter*{Bibliography}
\addcontentsline{toc}{chapter}{Bibliography}
\noindent The entries below are the sources already named at the ends of the chapters, in the order of those chapters. No source has been added.

\begin{enumerate}
\item J.~Bertrand, Solution d'un probl\`eme, \emph{Comptes Rendus de l'Acad\'emie des Sciences} \textbf{105} (1887), 369. The statement of the ballot probability.
\item D.~Andr\'e, Solution directe du probl\`eme r\'esolu par M.~Bertrand, \emph{Comptes Rendus de l'Acad\'emie des Sciences} \textbf{105} (1887), 436--437. A direct count of the favourable ballot orders.
\item W.~Feller, \emph{An Introduction to Probability Theory and Its Applications}, vol.~I, 3rd ed., Wiley, 1968, Chapter~III. The reflection comparison for paths, in the form used in Chapter~1.
\item M.~Renault, ``Lost (and Found) in Translation: Andr\'e's Actual Method and Its Application to the Generalized Ballot Problem,'' \emph{American Mathematical Monthly} \textbf{115} (2008), 358--363. The distinction between Andr\'e's original count and geometric reflection.
\item W.~Penney, Problem 95: Penney-Ante, \emph{Journal of Recreational Mathematics} \textbf{2} (1969), 241. The waiting game, posed for HHH against HTH.
\item S.~D.~Poisson, \emph{Recherches sur la probabilit\'e des jugements en mati\`ere criminelle et en mati\`ere civile; pr\'ec\'ed\'ees des r\`egles g\'en\'erales du calcul des probabilit\'es}, Bachelier, Paris, 1837. The law of rare events. Full title verified against the Biblioth\`eque nationale de France catalogue.
\item V.~Goncharov, Sur la distribution des cycles dans les permutations, \emph{C.\ R.\ (Doklady) Acad.\ Sci.\ URSS} \textbf{35} (1942), 267--269; Du domaine d'analyse combinatoire, \emph{Bull.\ Acad.\ Sci.\ URSS S\'er.\ Math.} \textbf{8} (1944), 3--48. The mean $H_n$ is Proposition~\ref{prop:cycles}. The Gaussian limit is Exercise~\ref{ex:cyclesclt} in Chapter~6.
\item W.~Feller, \emph{An Introduction to Probability Theory and Its Applications}, vol.~I, 3rd ed., Wiley, 1968. The Poisson limit, the coupon collector, and records.
\item F.~Galton, Problem 4001, \emph{Educational Times} \textbf{26} (1 April 1873), 17. The surname question. H.~W.~Watson’s solution is in the same journal, 1 August 1873.
\item H.~W.~Watson and F.~Galton, On the probability of the extinction of families, \emph{Journal of the Anthropological Institute} \textbf{4} (1875), 138--144. The functional equation in print.
\item W.~Feller, \emph{An Introduction to Probability Theory and Its Applications}, vol.~I, 3rd ed., Wiley, 1968, Chapter~XI for generating functions, and Chapter~XII for extinction and the total progeny.
\item J.~Bernoulli, \emph{Ars Conjectandi}, Basel, 1713. The weak law for frequencies, proved there by direct binomial estimates.
\item P.~L.~Chebyshev, Des valeurs moyennes, \emph{J.\ Math.\ Pures Appl.} (2) \textbf{12} (1867), 177--184. The inequality Volume~1 uses.
\item W.~Hoeffding, Probability inequalities for sums of bounded random variables, \emph{J.\ Amer.\ Statist.\ Assoc.} \textbf{58} (1963), 13--30. Theorem~\ref{thm:hoeffding}.
\item \'E.~Borel, Les probabilit\'es d\'enombrables et leurs applications arithm\'etiques, \emph{Rend.\ Circ.\ Mat.\ Palermo} \textbf{27} (1909), 247--271. The first lemma, for independent digits. F.~P.~Cantelli, Sulla probabilit\`a come limite della frequenza, \emph{Atti Accad.\ Naz.\ Lincei} \textbf{26} (1917), 39--45, removed the independence hypothesis from that half.
\item S.~Bernstein, D\'emonstration du th\'eor\`eme de Weierstrass fond\'ée sur le calcul des probabilit\'es, \emph{Comm.\ Soc.\ Math.\ Kharkov} \textbf{13} (1912), 1--2. Theorem~\ref{thm:weierstrass}.
\item G.~Vitali, Sul problema della misura dei gruppi di punti di una retta, Bologna, 1905. The obstruction to a length for every subset, under countable additivity and translation invariance.
\item Volume~1, Chapter~8, for the Riemann integral, the fundamental theorem, and $\int_0^{\infty}e^{-u^2}\,\dd u=\sqrt{\pi}/2$.
\item G.~E.~P.~Box and M.~E.~Muller, A note on the generation of random normal deviates, \emph{Ann.\ Math.\ Statist.} \textbf{29} (1958), 610--611. Theorem~\ref{thm:boxmuller}.
\item A.~de Moivre, \emph{Approximatio ad summam terminorum binomii}, privately printed in 1733 and reprinted in \emph{The Doctrine of Chances}, 1738. The local normal approximation for a fair binomial. For Stirling's formula, see Volume~1, Chapter~8.
\item J.~W.~Lindeberg, Eine neue Herleitung des Exponentialgesetzes in der Wahrscheinlichkeitsrechnung, \emph{Math.\ Z.} \textbf{15} (1922), 211--225. Theorem~\ref{thm:lindeberg}.
\item V.~Glivenko, Sulla determinazione empirica delle leggi di probabilit\`a, \emph{Giorn.\ Ist.\ Ital.\ Attuari} \textbf{4} (1933), 92--99. Theorem~\ref{thm:gc}, with Cantelli's note in the same volume, 421--424.
\item P.~L\'evy, \emph{Calcul des probabilit\'es}, Gauthier-Villars, 1925. The inversion formula in Theorem~\ref{thm:inversion} is the form used here; the continuity theorem in the converse direction is deferred to Volume~11.
\item C.~E.~Shannon, A Mathematical Theory of Communication, \emph{Bell System Technical Journal} \textbf{27} (1948), 379--423 and 623--656. Original entropy and source-coding framework.
\item D.~A.~Huffman, A Method for the Construction of Minimum-Redundancy Codes, \emph{Proceedings of the IRE} \textbf{40} (1952), 1098--1101. Original optimal merge rule for finite-alphabet prefix codes.
\item G.~P\'olya, \"Uber eine Aufgabe der Wahrscheinlichkeitsrechnung betreffend die Irrfahrt im Stra\ss ennetz, \emph{Mathematische Annalen} \textbf{84} (1921), 149--160, DOI \url{https://doi.org/10.1007/BF01458701}. The recurrence/transience theorem.
\item W.~Feller, \emph{An Introduction to Probability Theory and Its Applications}, vol.~I, 3rd ed., Wiley, 1968, Chapter~III. Reflection and ruin.
\item A.~A.~Markov, Extension of the law of large numbers to dependent quantities, \emph{Izvestiya Fiziko-matematicheskogo obshchestva pri Kazanskom universitete}, 2nd series, \textbf{15} (1906), 135--156.
\item W.~Doeblin, Expos\'e de la th\'eorie des cha\^ines simples constantes de Markoff \`a un nombre fini d'\'etats, \emph{Rev.\ Math.\ Union Interbalkan.} \textbf{2} (1938), 77--105.
\item D.~Aldous and P.~Diaconis, Shuffling Cards and Stopping Times, \emph{American Mathematical Monthly} \textbf{93} (1986), no.~5, 333--348. DOI: \texttt{10.1080/00029890.1986.11971821}. Theorem~\ref{thm:top}.
\item D.~Bayer and P.~Diaconis, Trailing the dovetail shuffle to its lair, \emph{Ann.\ Appl.\ Probab.} \textbf{2} (1992), 294--313. Theorem~\ref{thm:bayerdiaconis} and the distances after the statement of Proposition~\ref{prop:rifflecompose}. The three-decimal values for $52$ cards are the formula rounded, not a second model.
\item S.~D.~Poisson, \emph{Recherches sur la probabilit\'e des jugements}, Bachelier, Paris, 1837.
\item L.~Le Cam, An approximation theorem for the Poisson binomial distribution, \emph{Pacific J.\ Math.} \textbf{10} (1960), 1181--1197. Theorem~\ref{thm:lecam}, in the form proved in Chapter~10.
\item W.~Feller, \emph{An Introduction to Probability Theory and Its Applications}, Vol.~II, 2nd ed., Wiley, 1971. The comparison of a typical gap with the gap that covers a fixed time.
\item J.~L.~Doob, \emph{Stochastic Processes}, Wiley, 1953. Theorem~\ref{thm:doob} and Theorem~\ref{thm:L2conv} in the finite-partition form used here.
\item D.~Williams, \emph{Probability with Martingales}, Cambridge University Press, 1991. The casino proof of Theorem~\ref{thm:monkey}.
\item G.~P\'olya, Sur quelques points de la th\'eorie des probabilit\'es, \emph{Annales de l'Institut Henri Poincar\'e} \textbf{1} (1930), 117--161. The urn of Theorem~\ref{thm:polyaurn}.
\item N.~Wiener, Differential space, \emph{J.\ Math.\ and Phys.} \textbf{2} (1923), 131--174.
\item P.~L\'evy, \emph{Processus stochastiques et mouvement brownien}, Gauthier-Villars, 1948. The interpolation of Theorem~\ref{thm:levy}.
\item P.~Erd\H os, Some remarks on the theory of graphs, \emph{Bull.\ Amer.\ Math.\ Soc.} \textbf{53} (1947), 292--294. Theorem~\ref{thm:erdos}.
\item P.~Erd\H os and L.~Lov\'asz, Problems and results on $3$-chromatic hypergraphs and some related questions, in \emph{Infinite and Finite Sets}, North-Holland, 1975. Theorem~\ref{thm:lll}.
\item P.~Erd\H os and A.~R\'enyi, On the evolution of random graphs, \emph{Publ.\ Math.\ Inst.\ Hungar.\ Acad.\ Sci.} \textbf{5} (1960), 17--61. The component transition at $p=1/n$.
\item P.~Tur\'an, Eine Extremalaufgabe aus der Graphentheorie, \emph{Mat.\ Fiz.\ Lapok} \textbf{48} (1941), 436--452.
\item F.~P.~Ramsey, On a problem of formal logic, \emph{Proc.\ London Math.\ Soc.} \textbf{30} (1930), 264--286.
\item P.~Erd\H os, C.~Ko and R.~Rado, Intersection theorems for systems of finite sets, \emph{Quart.\ J.\ Math.\ Oxford} \textbf{12} (1961), 313--320.
\item G.~O.~H.~Katona, A simple proof of the Erd\H os--Ko--Rado theorem, \emph{J.\ Combin.\ Theory Ser.\ B} \textbf{13} (1972), 183--184.
\item P.~Erd\H os and R.~Rado, Intersection theorems for systems of sets, \emph{J.\ London Math.\ Soc.} \textbf{35} (1960), 85--90. The sunflower lemma.
\item B.~L.~van der Waerden, Beweis einer Baudetschen Vermutung, \emph{Nieuw Arch.\ Wisk.} \textbf{15} (1927), 212--216. Theorem~\ref{thm:vdw}. Theorem~\ref{thm:vdw3} is the case of two colours and progressions of length three, checked by hand.
\end{enumerate}
